\section{Colliding zeros, discriminants, and independence of the key invariant}
\label{sec:ftb-key-section-independence}
Retain the actual key datum and numerical parameters of
Section~\ref{sec:ftb-koszul-key}.  The section tuple is now allowed to
have nonreduced zeros in the special fibre.  The distinction matters:
the Jacobian factors need not be units individually.
The source cancels their common discriminant in the complete formula.

\subsection{The formal complete-intersection calculation}
\begin{srlemma}[The determinant class in the socle]
\label{lem:ftb-socle-determinant}
Let \(k\) be a field, \(S=k[[X_1,\ldots,X_g]]\), and
\(F_i=\sum_j a_{ij}X_j\).  Suppose every \(X_i^N\) belongs to
\((F_1,\ldots,F_g)\), for some \(N\).
Then \(\det(a_{ij})\notin(F_1,\ldots,F_g)\).
\end{srlemma}
\begin{proof}
Write \(X_i^{N+1}=\sum_jb_{ij}F_j\), put \(B=(b_{ij})\),
and let \(I=(X_i^{N+1}:1\leq i\leq g)\).
If \(\det a\in(F)\), the adjugate identity for \(B\), applied to
\(BF=(X_i^{N+1})\), gives \(\det(Ba)\in I\).
Also
\[
 (Ba)X=\operatorname{diag}(X_i^N)X .
\]
For a regular sequence \(X_1,\ldots,X_g\), two coefficient matrices
giving the same product with \(X\) have determinants equal modulo
the ideal of their common row products.  Here is the elementary
determinant argument used in the source.  Replace the rows one at
a time.  The difference row is a syzygy on the \(X_i\).
Expand the determinant difference against that row; the cofactors
\(C_j\) satisfy \(C_jX_l-C_lX_j\) in the common-product ideal,
by the adjugate identity with one row replaced by a coordinate row.
Induction on the support of a syzygy, using
\((X_j:j\in T):X_i=(X_j:j\in T)\) for \(i\notin T\), makes
its scalar product with these cofactors lie in that ideal.
Thus
\[
 \det(Ba)-\prod_iX_i^N\in I .
\]
It would follow that \(\prod_iX_i^N\in I\).
Every coefficient of a series in \(I\) whose exponent in every
variable is at most \(N\) is zero.  The coefficient of
\(\prod_iX_i^N\) in that monomial is one, a contradiction.
The colon identity above follows directly by inspecting coefficients
of \(pX_i\); it does not use an inverse to \(N!\).
\end{proof}

\begin{srlemma}[Discriminant and formal Jacobian]
\label{lem:ftb-formal-discriminant-jacobian}
Let \(R\) be local, \(S=R[[X_1,\ldots,X_g]]\), and
\(B=S/(F_1,\ldots,F_g)\) be free of finite rank over \(R\).
Assume the constant coefficients of the \(F_i\) lie in the maximal
ideal of \(R\).  For any \(R\)-basis \(b\) of \(B\),
\[
 \operatorname{disc}(b)\ \sim\
       N_{B/R}\!\left(\det(\partial_jF_i)\bmod(F)\right),
 \tag{K13}\label{eq:ftb-disc-jac}
\]
where \(\sim\) means association by an \(R\)-unit.
\end{srlemma}
\begin{proof}
The proof constructs the diagonal element used in the source.
Let
\(\mu:S\otimes_R B\to B\) be multiplication.
Its kernel is generated by
\(X_j\otimes1-1\otimes\bar X_j\).
To see this for power series, first write
\(s=C(s(0))+\sum_jX_js_j\).
If \(J\) is the ideal generated by those differences and
\(I\) is generated by \(X_j\otimes1\), this formula gives
\(\ker\mu\subseteq J+I\ker\mu\).
The kernel is finitely generated as an \(S\)-module:
\(B\) has a finite presentation over \(S\), and
\(S\otimes_R B\) is finite over \(S\).
The \(X_j\)'s are in the Jacobson radical, also after the finite
extension \(S\to S\otimes_RB\).
Nakayama therefore gives \(\ker\mu=J\).

Choose \(a_{ij}\in S\otimes_RB\) with
\[
 F_i\otimes1=\sum_j a_{ij}
              (X_j\otimes1-1\otimes\bar X_j).
\]
Map \(\det a\) to \(B\otimes_RB\), obtaining \(\Delta\).
The adjugate identity shows that
\[
       \Delta(z\otimes1-1\otimes z)=0\quad(z\in B).
 \tag{K14}\label{eq:ftb-diagonal-annihilator}
\]
Apply the derivation \(\partial_j\) in the first tensor factor and
then multiplication.  The Leibniz rule, with all diagonal
differences zero, gives
\(\mu(a_{ij})=\overline{\partial_jF_i}\).
Consequently the multiplication image of \(\Delta\) is
\[
                  \delta=\det(\partial_jF_i)\bmod(F).
\]

Write \(\Delta=\sum_{i,j}M_{ij}b_i\otimes b_j\).
We verify that \(M\) has unit determinant.
After reduction modulo the maximal ideal, the ideal generated by
the \(\bar X_j\)'s in \(B\) is nilpotent.  Indeed these elements
are in the radical of the base maximal ideal: finiteness over the
base and the Jacobson property of \(X_j\) give this radical
containment, and finite generation gives a common power.
The reduced quotient is thus supported at the origin.
Set the second tensor factor to its residue at the origin.
The image of the equations above is a matrix expression
\(\bar F_i=\sum_j\bar a_{ij}X_j\).
A common power of each \(X_j\) is in \((\bar F)\).
Lemma~\ref{lem:ftb-socle-determinant} says that the corresponding
image of \(\Delta\) is nonzero.

Here is the resulting matrix-unit argument.  Some row
\(w=\sum_jM_{i_0j}b_j\) has nonzero residue at the origin.
It is a unit of \(B\): it is a base scalar with nonzero residue
plus an element whose power lies in the base maximal ideal.
For any \(t=(t_j)\in R^{\operatorname{rank} B}\), put
\(y=w^{-1}\sum_jt_jb_j\).
Apply the linear functional \(b_{i_0}^*\) in the first factor of
\eqref{eq:ftb-diagonal-annihilator}, after multiplication by \(y\).
It gives
\[
 \sum_{i,j}M_{ij}b_{i_0}^*(b_i y)b_j=\sum_jt_jb_j .
\]
Thus right multiplication by \(M\) on row vectors is surjective.
A square matrix over a commutative ring with this property has unit
determinant.

Finally \eqref{eq:ftb-diagonal-annihilator}, with the same coordinate
functional and then summed over the basis index, gives
\[
 \sum_{i,j}M_{ij}\operatorname{Tr}_{B/R}(b_i z)b_j
                           =\delta z\quad(z\in B).
\]
If \(H_{ij}=\operatorname{Tr}_{B/R}(b_ib_j)\), the matrix of
multiplication by \(\delta\) is \(M^tH\).
Taking determinants yields
\(N_{B/R}(\delta)=\det(M)\det(H)\).
Since \(\det(M)\) is a unit, this is
\eqref{eq:ftb-disc-jac}.
\end{proof}

\begin{srlemma}[The order inside a split generic algebra]
\label{lem:ftb-formal-order-discriminant}
Let \(O\) be a complete DVR and
\(F=(F_1,\ldots,F_g)\subset O[[X_1,\ldots,X_g]]\).
Supply finitely many \(O\)-algebra maps \(e_a:O[[X]]\to O\)
with \(e_a(F_i)=0\) and \(e_a(X_i)\) in the maximal ideal.
Assume their common kernel is \((F)\), and the evaluation map
contains \(cO^d\) for a nonzero \(c\in O\).
Then \(B=O[[X]]/(F)\) has an \(O\)-basis \(\beta_1,\ldots,\beta_d\);
its point matrix \(P=(e_a(\beta_b))\) has nonzero determinant, and
\[
                 (\det P)^2\sim
                       \prod_a e_a(\det(\partial_jF_i)).
 \tag{K15}\label{eq:ftb-point-discriminant}
\]
\end{srlemma}
\begin{proof}
Evaluation embeds \(B\) in \(O^d\) and contains \(cO^d\).
Thus it is a finite torsion-free module of rank \(d\), hence free.
After passing to the fraction field the point matrix diagonalizes
all multiplication maps.
It follows that
\(\operatorname{Tr}(z)=\sum_ae_a(z)\) and
\(N(z)=\prod_ae_a(z)\).
In particular the trace Gram matrix is \(P^tP\).
The constant coefficients of \(F_i\) are in the maximal ideal,
because their value at any \(e_a\) is zero and the \(X_i\)'s
evaluate in that ideal.  If there are no points, both sides of
\eqref{eq:ftb-point-discriminant} are one.
Otherwise apply Lemma~\ref{lem:ftb-formal-discriminant-jacobian}
and the preceding trace and norm identities.
\end{proof}

\subsection{The actual global-to-formal order}
Let \(O\) be a complete DVR and let a tuple \(s_i\) on \(A/O\)
have finite common zero set on the closed fibre.
Suppose its generic zero scheme is split, reduced, and enumerated
by sections \(z_1,\ldots,z_d\), with nonzero generic Jacobians.
The sections \(z_a\) may coincide on the closed fibre.
Their distinct closed specializations form a finite set \(Y\).
Smoothness gives, at every member \(y\in Y\), the actual formal
coordinate chart \(O[[X_1,\ldots,X_g]]\to A\), obtained from an
etale coordinate neighbourhood and its completion.
Choose its coordinate origin to be one of the sections specializing
to \(y\).  Each \(z_a\) in this block is then evaluation at
\(X_i\) in the maximal ideal.  The local equations of \(s_i\)
are \(F_{y,i}\).

For a fixed \(l>gn\), let \(V_l=H^0(A,L^l)\).
The source first proves that the images of \(V_l\) generate
\[
                 B=\prod_{y\in Y}O[[X]]/(F_{y,i})
 \tag{K16}\label{eq:ftb-local-order}
\]
as an \(O\)-module; this step precedes any choice of an \(O\)-basis
of that product.  We give the two completion arguments.
On the algebraic closure of the residue field, approximate any
formal series to order \(q\) by a germ in the coordinate
neighbourhood.  The actual interpolation part of the
Cech/Koszul argument in Lemma~\ref{lem:ftb-koszul-unit}
produces a single section matching all these germs modulo the
equations and \((X)^q\).
Let \(\rho_1,\ldots,\rho_N\) be the images of a basis of \(V_l\).
The descending subspaces
\[
 \{v\in k^N:\ \sum_kv_k\rho_k\in(F)+(X)^q\}
\]
stabilize.  If they stabilize at \(q\), approximation at \(q+1\)
gives \((X)^q\subset(F)+(X)^{q+1}\).
Nakayama in \(k[[X]]\) gives \((X)^q\subset(F)\).
Thus the same \(\rho_k\)'s generate the exact residue quotient,
simultaneously in all blocks.

Lift this surjectivity coefficient by coefficient.
For precision the linear map has domain the product of the \(N\)
scalar coefficients and all coefficients of \(g\) multiplier series
in every block; it sends \((v,H_{yi})\) to the coefficients of
\(\sum_kv_k\rho_{y,k}+\sum_iH_{yi}F_{yi}\).
It is surjective modulo a uniformizer.
If the residue field was enlarged, apply a residue-field linear
projection to the coefficients first, and lift their projected
values to \(O\).
Successively correct modulo \(\pi,\pi^2,\ldots\).
Products of complete coefficient modules are complete, and the
target is separated; their coefficientwise limits give an exact
solution.  Consequently the map \(V_l\to B\) is surjective.

The integral kernel of ordinary point evaluation \(V_l\to O^d\)
is exactly \(\sum_i s_i H^0(A,L^{l-n})\).
Generic exactness first puts a nonzero multiple of a kernel vector
in this image.  The image is saturated: if \(\pi v=Du\), reduce
the degree-two part of \(u\) modulo \(\pi\).
Positive Koszul exactness in the residue field expresses it as
\(D\bar w\).  Lift \(w\); then \(u-Dw\) is divisible by \(\pi\)
in the relevant degree, so \(v\) is again a boundary.
Induct on the power of \(\pi\) in the first nonzero multiple.
Generic surjectivity also shows that the point-evaluation image
contains \(cO^d\) for one \(c\ne0\), by clearing denominators for
the finitely many coordinate vectors.

It follows that evaluation \(B\to O^d\) is injective and contains
\(cO^d\).  Apply
Lemma~\ref{lem:ftb-formal-order-discriminant} separately to the
blocks and combine their bases.
There are matrices \(E_l\) and \(P_l\) over \(O\) such that
\[
 \operatorname{ev}_l=P_lE_l,\qquad
 E_l:V_l\twoheadrightarrow O^d,\qquad \det P_l\ne0,\qquad
       (\det P_l)^2\sim J,                          \tag{K17}\label{eq:ftb-global-order-factor}
\]
where \(J\) is the product of the Jacobians computed in the given
Lie basis and line charts.
To compare formal and Lie Jacobians, differentiate the actual
formal coordinate chart using the square-zero algebra
\(O\oplus O^g\).
The matrix of the translated Lie basis in the formal coordinate
tangents has a right inverse, by the Lie basis property and the
formal lifting property.  It therefore has unit determinant.
The chain rule multiplies each formal Jacobian by this unit.
Changing the local line generator adds a derivative of that generator
times \(s_i(z_a)=0\), so it only contributes another unit.
This proves precisely the association in
\eqref{eq:ftb-global-order-factor}, with no assumption of distinct
closed specializations.

Replacing the augmentation \(\operatorname{ev}_l\) of the Koszul
complex by \(E_l\) gives a complex with exact residue fibre.
Indeed \(E_l\) is surjective.  A residue vector in its kernel lifts
to an integral vector; subtract integral preimages of its evaluation
divided by \(\pi\), using a fixed right inverse to \(E_l\).
Its residue is unchanged and its exact evaluation is zero.
The saturated kernel calculation above then supplies the
degree-two preimage.  In higher degrees use the same residue
Koszul exactness as before.  Thus its squared torsion \(t_l^0\)
is an \(O\)-unit.  Changing the augmentation from \(E_l\) to
\(P_lE_l\) multiplies squared torsion by \((\det P_l)^2\).
This is the source finite-closed-fibre factorization, including
the integral order, not only the generic reduced zero scheme.

\subsection{Unit cancellation for the full invariant}
\begin{srproposition}[Finite closed-fibre zeros suffice]
\label{prop:ftb-key-finite-zeros-unit}
Let \(O\) be a DVR containing \(\mathbf Q\), and let \(x\) be an
actual key datum over \(O\).  Let \(s_i\) have a finite common
zero set on the closed fibre.  Let \(F\) be any field with an
injective map \(O\to F\), carrying a base change of \(x,s_i\)
and an actual value \(\kappa\) of
\eqref{eq:ftb-kaform} on a split transverse generic intersection.
If \(a\in O\), \(e\in\mathbf Z\setminus\{0\}\), and
\(\kappa^e=a\) in \(F\), then \(a\in O^\times\).
\end{srproposition}
\begin{proof}
Only finitely many points occur in the generic zero enumeration and
the two systems of split kernel tuples.
Descend their coordinates to a finitely generated subfield of \(F\)
over \(\operatorname{Frac}O\), choose a discrete valuation ring
dominating \(O\) in that field, and complete it.
The source dominating-DVR construction applies to an essentially
finite type field extension of a nonfield DVR.
Properness of \(A/O\) extends all the descended points to this
valuation ring.  A common overfield compares its completion with
\(F\); take the compositum through a maximal ideal of the tensor
product of the two fields.
This gives a complete DVR \(O'\), a local injective map \(O\to O'\),
and an injective map \(O'\to\Omega\), with all required points
integral and the original generic data preserved.
The closed-fibre zero set remains finite by quasi-finiteness of
the reduced finite closed subscheme and base change.
Because \(O\) contains \(\mathbf Q\), all positive integers,
in particular every auxiliary \(D\), remain units.

Use \eqref{eq:ftb-global-order-factor} for \(l\) and \(l'\).
After accounting for the generic line-chart factors, put
\(\gamma=\prod_a g_a\).
There are units \(t,t',\eta,\eta'\in O'^\times\) and nonzero
generic elements \(\delta,\delta'\) such that the two generic
squared Koszul torsions and the generic Jacobian product are
\[
 \begin{aligned}
 t_l&=t\,\gamma^{-2l}\delta^2,&
 t_{l'}&=t'\,\gamma^{-2l'}(\delta')^2,\\
 J_{\rm gen}&=\eta\,\gamma^{-gn}\delta^2
            =\eta'\,\gamma^{-gn}(\delta')^2 .
 \end{aligned}                                      \tag{K18}\label{eq:ftb-colliding-factors}
\]
Here \(\delta=\det P_l\) and \(\delta'=\det P_{l'}\).
The cubical determinant units \(\lambda_j,\lambda'_j\) are integral
units by Lemma~\ref{lem:ftb-box-det-unit}; the prime-to-residue
kernel tuples are fibrewise distinct, since their multiplication
maps are etale and \(O'\) is connected.

Now \((\delta')^2=(\eta/\eta')\delta^2\).
Substitute \eqref{eq:ftb-colliding-factors} into
\eqref{eq:ftb-kaform}.
The exponent of \(\gamma\) is the zero expression already computed
in Proposition~\ref{prop:ftb-integral-koszul-invariant}.
The exponent of the remaining \(\delta^2\) is
\[
 4D^{3g+2}\{(gn-2l')+(2l-gn)\}
                    +8D^{3g+2}(l'-l)=0 .
\]
Everything left is an \(O'\)-unit.  Uniqueness for the fixed generic
intersection identifies this unit with the original \(\kappa\).
Its nonzero integer power \(a\) is therefore a unit in \(O'\).
A local map reflects units, so \(a\) is a unit in \(O\).
\end{proof}

\subsection{Constancy along lines and passage to all tuples}
\begin{srlemma}[Source line argument]\label{lem:ftb-key-line}
Let \(K\) be algebraically closed of characteristic zero.
Suppose two transverse tuples \(s_0,s_1\) on a fixed key datum have
the property that the common zero set of
\((1-t)s_0+ts_1\) is finite for every \(t\in K\).
Then their values in \eqref{eq:ftb-kaform} are equal.
\end{srlemma}
\begin{proof}
Work over the parameter line \(K[T]\).
The transverse data at \(T=a\) lift after a local DVR extension of
\(K[T]_{(T-a)}\).  Indeed take a complete DVR whose residue field
contains the finitely many coordinates of the residue zeros and
prime-to-\(D\) torsion.  The Jacobian criterion makes each zero
etale; Hensel lifting gives its unique section and unit Jacobian.
Properness ensures every common zero specializes to a listed
residue zero; the local etale uniqueness therefore lists all zeros.
The split etale \(D\)-torsion lifts as well.
Proposition~\ref{prop:ftb-integral-koszul-invariant} constructs an
integral unit value there, specializing to the given one.

Field-base-change compatibility and uniqueness descend this generic
value to \(K(T)\).  Explicitly, two embeddings of an extension field
over \(K(T)\) give the same value by uniqueness.  The element is
therefore purely inseparable over \(K(T)\): its simple intermediate
field has at most one embedding into any common field.
In characteristic zero it belongs to \(K(T)\) itself.
Call it \(c(T)\ne0\).
At every \(T=a\), Proposition~\ref{prop:ftb-key-finite-zeros-unit}
says that a nonzero integer power of \(c(T)\) lying in the local
DVR is a unit.  Write \(c=P/Q\) in relatively prime polynomials.
At a root of \(P\), \(Q\) is nonzero and this criterion applied with
exponent \(1\) is a contradiction.  At a root of \(Q\), use
exponent \(-1\).
Thus \(P,Q\) have no roots in \(K\), so they are constants.
Specialization of the lifted unit then identifies both endpoint
values with this same constant.
\end{proof}

For use in this line argument, the source proves the following
finiteness statement.  Over an algebraically closed field \(k\),
let \(s\) have finite common zeros, and replace one section by the
line towards a section whose coordinates \(T_j\) in a basis of
\(H^0(L^n)\), \(n\geq4\), are algebraically independent over \(k\).
After any algebraically closed extension, every member of that
line still has finite common zeros.
Here is the proof.  Removing one equation gives a closed set \(X\)
of dimension at most one: an ample divisor on a positive-dimensional
proper component has nonempty intersection, and the principal-ideal
dimension inequality bounds that component by one when its
intersection with the original equation is finite.
If an altered zero set were infinite, a curve component would
supply two distinct closed \(k\)-points \(y_1,y_2\) outside the
original finite zero set, whose extension fibres lie in it.
(The curve is geometrically integral after selecting an irreducible
component over algebraically closed \(k\).)
Let \(a_\nu\) be the original section values and
\(\beta_{\nu j}\) the basis-section values, using line generators
at the two points.  The line equations give
\[
 (1-t)a_\nu+t\sum_jT_j\beta_{\nu j}=0 .
\]
Take \(y_1\) outside the original zero set, so \(a_1\ne0\).
Then \(t\ne0\); eliminating \(t\) and using linear independence
of the \(T_j\)'s gives
\(a_2\beta_{1j}=a_1\beta_{2j}\) for every \(j\).
Every section vanishing at \(y_1\) must therefore vanish at \(y_2\).
This contradicts point separation by \(H^0(L^n)\).
The argument permits arbitrary \(t\) in the extension field.

\begin{srproposition}[Independence of the auxiliary sections]
\label{prop:ftb-key-tuple-independent}
Fix admissible \(g,n,l,l',\alpha\) in characteristic zero.
For each key datum over an algebraically closed field \(K\), there is
one value \(\kappa_{n,l,l',\alpha}(x)\in K^\times\) such that
every actual transverse-intersection construction
\eqref{eq:ftb-kaform}, after any algebraically closed field
extension, gives its image.  This value exists on an extension,
is functorial in the framed key datum, and is independent of the
auxiliary section tuple.
\end{srproposition}
\begin{proof}
First take two coordinate tuples jointly algebraically independent
over a field of definition for \(x\) and a basis of \(H^0(L^n)\).
Generic sections meet transversally.  To verify the latter source
step, at a proposed zero descend its values and first jets to a field
of transcendence degree at most \(g\), using the dimension of \(A\).
Evaluation is one nonzero linear equation on each coefficient row.
Eliminate one coefficient per row.  The remaining determinant of
the first jets is a nonzero polynomial: first-jet separation supplies
one solution of the evaluation equations with nonzero determinant.
Its evaluation at the remaining generic coefficients is injective.
Indeed its polynomial ring has transcendence degree at most the
number of original independent coefficients, while those
coefficients belong to the image; a nonzero kernel would lower
that degree.  Thus the generic Jacobian is nonzero at every zero.
Smoothness makes these zeros isolated; properness makes their set
finite, and they are rational over the algebraically closed field.
This is the source transcendence-degree proof of generic
transversality.

Replace rows of the first generic tuple by rows of the second, one
at a time.  For each replacement, pass to the algebraic closure of
the field generated by all the retained rows.
The incoming row remains algebraically independent over this field.
The preceding finiteness argument and
Lemma~\ref{lem:ftb-key-line} give equal values at successive tuples.
Thus jointly generic tuples give the same value.
For any two separately generic tuples, adjoin a third independent
tuple; comparison through it gives the same conclusion.

Choose the universal generic tuple in
\(\Omega=\overline{K(T_j)}\).
Every \(K\)-embedding \(\Omega\to\Omega\) takes it to another
generic tuple, so fixes its invariant.
Its invariant therefore lies in \(K\).
For clarity, a nonconstant transcendental element would extend to a
transcendence basis and could be moved by the embedding replacing
each basis variable by that variable plus one; an element algebraic
over algebraically closed \(K\) is already in \(K\).

Finally let \(S_j\) be any tuple over a field extension for which an
actual transverse value is given.
Over
\(\Omega'=\overline{F(U_j)}\), use the complete DVR
\(\Omega'[[t]]\) and coordinates
\(S_j+(U_j-S_j)t\).
They specialize to \(S_j\), and are algebraically independent over
\(K\), since evaluation at \(t=1\) gives the independent \(U_j\)'s.
Lift the transverse special-fibre data by the Hensel construction
in Lemma~\ref{lem:ftb-key-line}.
Its generic value is the universal generic value; specializing the
integral unit identifies the original value with it.
Base change and transport of the framed datum follow by applying
the already proved matrix base-change law and fixed-intersection
uniqueness.  Nonzeroness holds because the generic construction is
a product of units.
\end{proof}

\paragraph{Remaining key-rule scope.}
The argument above discharges section independence and the
finite-closed-fibre unit calculation, relative only to the specified
abelian-scheme, smooth formal-coordinate, and dominating-DVR laws.
It supplies an integral algebraic invariant with the positive
exponent \(c=\operatorname{kaExp}(g,n,l,l')\).
The source still passes from this auxiliary invariant to its
normalized analytic value, replaces the auxiliary residue prime to
get all-place units, and uses connected families to obtain the
exponent-one \(\kappa\).
Those further analytic and family statements are not consequences
of section independence alone.
